\documentclass[11pt]{article}
\usepackage{amsmath,amssymb,amsthm}

\newtheorem{theorem}{Theorem}
\newtheorem{corollary}[theorem]{Corollary}

\newcommand{\Up}{U_+}
\newcommand{\mass}{\operatorname{B}}

\title{The complete excess-three inverse decoder layer}
\author{}
\date{27 August 2026}

\begin{document}
\raggedbottom
\maketitle

For a positive odd integer $x$, let
\[
 \Up(x):=\frac{3x+1}{2^{v_2(3x+1)}}.
 \tag{1}
\]
An exponent word $E=(e_1,\ldots,e_h)$ is read from its positive odd
start toward its terminal state.  We define its mass and excess by
\[
 \mass(E):=e_1+\cdots+e_h,
 \qquad t(E):=\mass(E)-h.
 \tag{2}
\]
The purpose of this note is to classify every exact inverse word with
$t(E)=3$ at a fixed positive odd terminal state.

We begin by recalling the reverse rule.  If $x$ is the current odd
terminal state, an exponent-$a$ predecessor is
\[
 p=\frac{2^a x-1}{3}.
 \tag{3}
\]
It exists precisely when $3\nmid x$ and
\[
 a\equiv
 \begin{cases}
 1\pmod2,&x\equiv2\pmod3,\\
 0\pmod2,&x\equiv1\pmod3.
 \end{cases}
 \tag{4}
\]
In the shifted coordinate $X:=x+1$, the four reverse operations used
below are
\[
 X\longmapsto\frac{2X}{3},
 \qquad
 X\longmapsto\frac{4X-2}{3},
 \qquad
 X\longmapsto\frac{8X-6}{3},
 \qquad
 X\longmapsto\frac{16X-14}{3}.
 \tag{5}
\]

\begin{theorem}[complete excess-three classification]
Let $r$ be a positive odd integer and put
\[
 b:=v_3(r+1),
 \qquad s:=\frac{r+1}{3^b}.
 \tag{6}
\]
There are exactly four families of exact predecessor words $E$ with
$t(E)=3$.

\medskip
\noindent\textbf{Family I: one exponent four.}
Define
\[
 Q:=2^{b+3}s-7,
 \qquad c:=v_3(Q).
 \tag{7}
\]
For every $1\leq j\leq c$, there is one exact word
\[
 \boxed{
 E=1^{j-1}4\,1^b,
 \qquad |E|=b+j,
 \qquad z+1=\left(\frac23\right)^jQ.}
 \tag{8}
\]
This family is empty when $c=0$.

\medskip
\noindent\textbf{Family II: an exponent three followed in reverse by an
exponent two.}
For every $0\leq k<b$, define
\[
 M_k:=2^{k+2}3^{b-k-1}s-1,
 \qquad c_k:=v_3(M_k),
 \qquad u_k:=\frac{M_k}{3^{c_k}},
 \tag{9}
\]
We also define
\[
 P_k:=2^{c_k+2}u_k-1,
 \qquad d_k:=v_3(P_k).
 \tag{10}
\]
For every $1\leq j\leq d_k$, there is one exact word
\[
 \boxed{
 E=1^{j-1}2\,1^{c_k}3\,1^k,
 \qquad |E|=k+c_k+j+1,
 \qquad z+1=\left(\frac23\right)^jP_k.}
 \tag{11}
\]
The branch indexed by $k$ is empty when $d_k=0$.

\medskip
\noindent\textbf{Family III: an exponent two followed in reverse by an
exponent three.}
Define
\[
 N_1:=2^{b+1}s-1,
 \qquad b_1:=v_3(N_1).
 \tag{12}
\]
This family is empty when $b_1=0$.  If $b_1\geq1$, put
\[
 s_1:=\frac{N_1}{3^{b_1}}.
 \tag{13}
\]
For every $0\leq k<b_1-1$, define
\[
 R_k:=2^{k+3}3^{b_1-k-2}s_1-1,
 \qquad f_k:=v_3(R_k).
 \tag{14}
\]
For every $0\leq\ell\leq f_k$, there is one exact word
\[
 \boxed{
 E=1^\ell3\,1^k2\,1^b,
 \qquad |E|=b+k+\ell+2,
 \qquad z+1=2\left(\frac23\right)^\ell R_k.}
 \tag{15}
\]

\newpage
\noindent\textbf{Family IV: three exponents two.}
Suppose $b_1\geq1$, and use $s_1$ from (13).  Define
\[
 N_2:=2^{b_1+1}s_1-1,
 \qquad b_2:=v_3(N_2).
 \tag{16}
\]
This family is empty when $b_2=0$.  If $b_2\geq1$, put
\[
 s_2:=\frac{N_2}{3^{b_2}},
 \qquad
 N_3:=2^{b_2+1}s_2-1,
 \qquad b_3:=v_3(N_3).
 \tag{17}
\]
For every $1\leq j\leq b_3$, there is one exact word
\[
 \boxed{
 \begin{aligned}
 E&=1^{j-1}2\,1^{b_2-1}2\,1^{b_1-1}2\,1^b,\\
 |E|&=b+b_1+b_2+j,\\
 z+1&=\left(\frac23\right)^jN_3.
 \end{aligned}}
 \tag{18}
\]
There are no other exact words of excess three.

For the horizon statement below, we set $b_2=b_3:=0$ when $b_1=0$, and
we set $b_3:=0$ when $b_1\geq1$ and $b_2=0$.
\end{theorem}

\begin{proof}
Every exponent contributes $e_i-1$ to the excess.  Therefore, an
excess-three word has exactly one of the following forms: one exponent
four, one exponent three together with one exponent two, or three
exponents two.  We read each possibility backward from $r$.

The shifted terminal state is $3^bs$.  Precisely $b$ consecutive
exponent-one inverses are initially available, and after all of them the
shifted state is $2^bs$.

Suppose first that the word contains one exponent four.  This exponent
must occur at the first even-parity junction, after all $b$
exponent-one inverses.  Reversing it gives
\[
 \frac{16\cdot2^bs-14}{3}=\frac{2Q}{3}.
 \tag{19}
\]
This is integral exactly when $c\geq1$.  Its $3$-adic valuation is
$c-1$, so exactly $j-1$ additional exponent-one inverses are available
for $1\leq j\leq c$.  Their shifted starts are $(2/3)^jQ$.  This proves
Family I.

Next suppose that an exponent three is encountered first when reading
backward.  It must occur strictly inside the initial exponent-one run.
After $k$ exponent-one inverses, where $0\leq k<b$, the shifted state is
\[
 2^k3^{b-k}s.
 \tag{20}
\]
Reversing exponent three gives $2M_k$.  The next non-one inverse must be
the unique exponent two, so all $c_k$ available exponent-one inverses
must first be taken.  The shifted state at that junction is
$2^{c_k+1}u_k$, and the exponent-two inverse gives
\[
 \frac{4\cdot2^{c_k+1}u_k-2}{3}=\frac{2P_k}{3}.
 \tag{21}
\]
The remaining exponent-one inverses give exactly (11).  This proves
Family II.

Suppose instead that the exponent two is encountered first.  It must
occur at the first even-parity junction.  Its shifted output is
\[
 \frac{4\cdot2^bs-2}{3}=\frac{2N_1}{3},
 \tag{22}
\]
so it exists exactly when $b_1\geq1$.  If the exponent three is
encountered next, it must occur strictly inside the following run of
$b_1-1$ exponent-one inverses.  After $k$ of these inverses the shifted
state is
\[
 2^{k+1}3^{b_1-k-1}s_1.
 \tag{23}
\]
Reversing exponent three gives $2R_k$, after which precisely $f_k$
exponent-one inverses remain available.  This proves Family III.

It remains to consider three exponents two.  Parity forces each one to
occur at the next junction, after the current exponent-one run is
exhausted.  The first junction is described by (22).  After the next
$b_1-1$ exponent-one inverses, the shifted state is $2^{b_1}s_1$, and
the second exponent-two inverse gives $2N_2/3$.  After the following
$b_2-1$ exponent-one inverses, the shifted state is $2^{b_2}s_2$, and
the third exponent-two inverse gives $2N_3/3$.  Reversing a final
$j-1$ exponent-one steps gives (18).

The three partitions of the excess are exhaustive.  In each partition,
the parity rule forces every non-one exponent into the positions used
above.  Therefore, the four displayed families contain every exact
excess-three word.
\end{proof}

\begin{corollary}[finite excess-three horizon]
With the notation of Theorem~1, define
\[
 H_I:=b+c,
 \tag{24}
\]
\[
 H_{II}:=\max_{0\leq k<b}(k+c_k+d_k+1),
 \qquad
 H_{III}:=\max_{0\leq k<b_1-1}(b+k+f_k+2),
 \tag{25}
\]
where the maximum of an empty set is zero, and define
\[
 H_{IV}:=
 \begin{cases}
 b+b_1+b_2+b_3,&b_1b_2b_3>0,\\
 0,&b_1b_2b_3=0.
 \end{cases}
 \tag{26}
\]
Then no exact $h$-step predecessor word of $r$ has excess three when
\[
 \boxed{h>H_3(r):=\max(H_I,H_{II},H_{III},H_{IV}).}
 \tag{27}
\]
If $H_{\leq2}(r)$ denotes the horizon from the excess-zero, excess-one,
and excess-two classifications, then every exact prefix with
\[
 h>H_{\leq3}(r):=\max\bigl(H_{\leq2}(r),H_3(r)\bigr)
 \tag{28}
\]
has excess at least four.
\end{corollary}

\begin{proof}
The largest word length in each family is the corresponding quantity in
(24)--(26).  Theorem~1 proves (27).  Combining it with the three earlier
decoder layers proves (28).
\end{proof}

\begin{corollary}[fixed-seed inequalities]
Let $n$ be the globally least positive odd integer with an unbounded
accelerated orbit, and suppose that the orbit of $r$ is unbounded.  Every
branch in Theorem~1 satisfies its corresponding inequality
\[
 \boxed{
 \begin{aligned}
 Q&\geq(n+1)\left(\frac32\right)^j,
 &&1\leq j\leq c,\\
 P_k&\geq(n+1)\left(\frac32\right)^j,
 &&1\leq j\leq d_k,\\
 R_k&\geq\frac{n+1}{2}\left(\frac32\right)^\ell,
 &&0\leq\ell\leq f_k,\\
 N_3&\geq(n+1)\left(\frac32\right)^j,
 &&1\leq j\leq b_3.
 \end{aligned}}
 \tag{29}
\]
In particular, every minimum-channel prefix of excess at most three is
now contained in a finite explicit valuation ladder.  Beyond
$H_{\leq3}(r)$, that channel has excess at least four.
\end{corollary}

\begin{proof}
Every positive predecessor of $r$ also has an unbounded forward orbit.
Its start is therefore at least $n$.  Substitution of the four shifted
starts in (8), (11), (15), and (18) proves (29).  The final assertion is
Corollary~2.
\end{proof}

The accompanying exact certificate is
\[
 \begin{gathered}
 \texttt{verify\_divergence\_record\_}\\
 \texttt{excess\_three.py}.
 \end{gathered}
\]
It independently enumerates every reverse word of excess three and checks the four
families for all $50{,}001$ odd terminal states $r\leq100{,}001$ using
exact integer arithmetic.

\end{document}
